% Problemas del capítulo 5. Generado por tools/problemas/gen.py; no editar a mano.
\section{Problemas}\label{sec:problemas-05}

Graduados: \(\star\) directo, \(\star\star\) para pensar, \(\star\star\star\) desafío. Las respuestas están en el apéndice~\ref{ap:respuestas}; conviene intentarlos antes de mirarlas.

\begin{problema}[\(\star\)]\label{pb:05-1}
Para la forma \(f(x,y)=3x-2y\), calculá \(f(4,1)\). Dibujá las rectas \(f=0\), \(f=1\) y \(f=2\): ¿qué distancia hay entre dos consecutivas?
\end{problema}

\begin{problema}[\(\star\star\)]\label{pb:05-2}
Un terreno plano tiene cota \(z=100+0{,}02\,E-0{,}01\,N\) (en metros). ¿Cuánto subís si caminás \(50\) m al este y \(30\) m al norte? ¿Qué forma lineal da esa subida?
\end{problema}

\begin{problema}[\(\star\star\)]\label{pb:05-3}
Repetí la cuenta del problema anterior para el paso \((3,4)\) m, pero midiendo los pasos en centímetros. ¿Cómo cambian las componentes del paso y las de la forma? ¿Cambia la subida?
\end{problema}

\begin{problema}[\(\star\star\)]\label{pb:05-4}
Encontrá la base dual de \(b_1=(1,1)\), \(b_2=(1,-1)\): las dos formas \(\beta_1\), \(\beta_2\) que valen \(1\) en su vector y \(0\) en el otro.
\end{problema}

\begin{problema}[\(\star\star\star\)]\label{pb:05-5}
En el terreno del problema 2, ¿hacia qué acimut se sube más rápido? ¿Cuál es esa pendiente máxima, en porcentaje?
\end{problema}
